\section{Семинар 1. Математическая индукция и битовые операции}
\noeqref{eq:fact_sum, eq:strong_ind, eq:x_inv_n, eq:x_inv_rec, eq:cubes_sum}
\noeqref{eq:hanoi_rec, eq:am_gm, eq:bin_rep, eq:ops_count, eq:cyc_shift_sum}

\begin{remark}[Обзор темы]
    Одним из основных инструментов доказательства утверждений, зависящих от 
    натурального параметра, является метод математической индукции. На этом 
    семинаре рассматриваются классическая, сильная и прямо-обратная 
    индукции, а также применение индуктивных рассуждений к алгоритмическим 
    задачам и анализу двоичного представления чисел.
\end{remark}

\subsection{Метод математической индукции (ММИ)}

\begin{definition}
    Пусть имеется некоторое утверждение $P(n)$, зависящее от натурального 
    параметра $n \in \N$. \textdef{Принцип математической индукции} гласит, 
    что утверждение $P(n)$ истинно для всех натуральных чисел $n \in \N$, 
    если одновременно выполняются два условия:
    \begin{enumerate}
        \item \textdef{База индукции:} утверждение $P(1)$ истинно;
        \item \textdef{Индукционный переход (шаг):} для любого натурального 
        $k \ge 1$ из допущения об истинности $P(k)$ (\textdef{индукционное 
        предположение}) логически следует истинность $P(k + 1)$.
    \end{enumerate}
\end{definition}

\begin{remark}
    Метод математической индукции строго обосновывается аксиоматикой Пеано 
    для натуральных чисел. В частности, пятая аксиома (аксиома индукции) 
    постулирует, что если некоторое подмножество $M \subseteq \N$ содержит 
    единицу ($1 \in M$) и вместе с каждым элементом $k$ содержит 
    следующий за ним элемент ($k \in M \implies k + 1 \in M$), то множество 
    $M$ совпадает со всем множеством натуральных чисел ($M = \N$).
\end{remark}

\begin{problem}
    Доказать тождество для суммы с факториалами при любом $n \in \N$:
    \begin{equation} \label{eq:fact_sum}
        \sum_{i=1}^n i \cdot i! = (n + 1)! - 1 \tag{1}
    \end{equation}
\end{problem}

\begin{proof}[Решение]~\\
    \begin{quote}
        Проведём доказательство тождества \eqref{eq:fact_sum} методом 
        математической индукции по параметру $n$.

        \textbf{1. База индукции ($n = 1$).} Вычислим левую и правую части 
        равенства \eqref{eq:fact_sum} при $n = 1$:
        \[
            \sum_{i=1}^1 i \cdot i! = 1 \cdot 1! = 1, \qquad
            (1 + 1)! - 1 = 2! - 1 = 2 - 1 = 1.
        \]
        Так как $1 = 1$, база индукции выполняется.

        \textbf{2. Индукционный переход ($k \implies k + 1$).} 
        Предположим, что формула \eqref{eq:fact_sum} верна для некоторого 
        произвольного натурального $k \ge 1$ (индукционное предположение):
        \[
            \sum_{i=1}^k i \cdot i! = (k + 1)! - 1
        \]
        Докажем справедливость формулы для $n = k + 1$, то есть покажем, что:
        \[
            \sum_{i=1}^{k+1} i \cdot i! = (k + 2)! - 1
        \]
        Выделим из суммы для $k + 1$ слагаемых последнее слагаемое и 
        применим индукционное предположение к сумме первых $k$ членов:
        \begin{align*}
            \sum_{i=1}^{k+1} i \cdot i! &= 
            \left( \sum_{i=1}^k i \cdot i! \right) + (k + 1) \cdot (k + 1)! \\
            &= \underbrace{\bigl((k + 1)! - 1\bigr)}_{\text{по предположению}} 
            + (k + 1) \cdot (k + 1)!
        \end{align*}
        Вынесем общий множитель $(k + 1)!$ за скобки:
        \begin{align*}
            (k + 1)! + (k + 1) \cdot (k + 1)! - 1 &= 
            (k + 1)! \cdot \bigl(1 + (k + 1)\bigr) - 1 \\
            &= (k + 1)! \cdot (k + 2) - 1 = (k + 2)! - 1.
        \end{align*}
        Получено в точности требуемое выражение для $n = k + 1$. 
        Индукционный переход доказан, следовательно, по принципу 
        математической индукции тождество \eqref{eq:fact_sum} верно для 
        всех $n \in \N$.
    \end{quote}
\end{proof}

\subsection{Сильная (возвратная) индукция}

\begin{definition}
    \textdef{Сильная (совершенная или возвратная) математическая индукция} 
    отличается от классической тем, что в индукционном переходе 
    предполагается истинность утверждения не только для непосредственно 
    предшествующего номера $k$, но и для \emph{всех} предыдущих номеров 
    от $1$ до $k$ включительно. Шаг индукции имеет вид:
    \begin{equation} \label{eq:strong_ind}
        \bigl(P(1) \land P(2) \land \dots \land P(k)\bigr) \implies P(k + 1) 
        \tag{2}
    \end{equation}
\end{definition}

\begin{problem}
    Пусть $x \in \R \setminus \{0\}$ --- некоторое вещественное число, 
    такое что $\left(x + \frac{1}{x}\right) \in \Z$. Доказать, что для 
    любого натурального $n \in \N$ выполняется соотношение:
    \begin{equation} \label{eq:x_inv_n}
        \left(x^n + \frac{1}{x^n}\right) \in \Z \tag{3}
    \end{equation}
\end{problem}

\begin{proof}[Решение]~\\
    \begin{quote}
        Докажем утверждение \eqref{eq:x_inv_n} методом сильной 
        математической индукции по $n$.

        \textbf{1. База индукции ($n = 1$ и $n = 2$).} 
        Поскольку на шаге индукции нам потребуется опираться на два 
        предыдущих значения параметра, проверим базу сразу для двух 
        начальных значений:
        \begin{itemize}
            \item При $n = 1$ условие $\left(x^1 + \frac{1}{x^1}\right) \in \Z$ 
            выполнено по условию задачи.
            \item При $n = 2$ возведём исходное выражение в квадрат:
            \[
                \left(x + \frac{1}{x}\right)^2 = 
                x^2 + 2 \cdot x \cdot \frac{1}{x} + \frac{1}{x^2} = 
                x^2 + 2 + \frac{1}{x^2}
            \]
            Выразим отсюда сумму квадратов:
            \[
                x^2 + \frac{1}{x^2} = \left(x + \frac{1}{x}\right)^2 - 2
            \]
            Так как $\left(x + \frac{1}{x}\right) \in \Z$, то его квадрат 
            и разность с числом $2$ также принадлежат $\Z$. База для 
            $n = 1, 2$ верна.
        \end{itemize}

        \textbf{2. Индукционный переход.} Пусть $k \ge 2$ и утверждение 
        \eqref{eq:x_inv_n} верно для всех натуральных чисел $m \le k$. 
        В частности, по предположению сильной индукции \eqref{eq:strong_ind} 
        числа $\left(x^k + \frac{1}{x^k}\right)$ и 
        $\left(x^{k-1} + \frac{1}{x^{k-1}}\right)$ являются целыми.

        Докажем, что $\left(x^{k+1} + \frac{1}{x^{k+1}}\right) \in \Z$. 
        Рассмотрим произведение:
        \[
            \left(x^k + \frac{1}{x^k}\right) \left(x + \frac{1}{x}\right) = 
            x^{k+1} + x^{k-1} + \frac{1}{x^{k-1}} + \frac{1}{x^{k+1}}
        \]
        Выразим искомую сумму для показателя $k + 1$:
        \begin{equation} \label{eq:x_inv_rec}
            x^{k+1} + \frac{1}{x^{k+1}} = 
            \left(x^k + \frac{1}{x^k}\right) \left(x + \frac{1}{x}\right) - 
            \left(x^{k-1} + \frac{1}{x^{k-1}}\right) \tag{4}
        \end{equation}
        В правой части тождества \eqref{eq:x_inv_rec}:
        \begin{itemize}
            \item $\left(x^k + \frac{1}{x^k}\right) \in \Z$ по 
            индукционному предположению для $m = k$;
            \item $\left(x + \frac{1}{x}\right) \in \Z$ по условию (база 
            $m = 1$);
            \item $\left(x^{k-1} + \frac{1}{x^{k-1}}\right) \in \Z$ по 
            индукционному предположению для $m = k - 1$.
        \end{itemize}
        Поскольку кольцо целых чисел $\Z$ замкнуто относительно операций 
        умножения и вычитания, правая часть \eqref{eq:x_inv_rec} является 
        целым числом. Следовательно, 
        $\left(x^{k+1} + \frac{1}{x^{k+1}}\right) \in \Z$, и по принципу 
        сильной индукции утверждение доказано для всех $n \in \N$.
    \end{quote}
\end{proof}

\subsection{Индукция в задачах на делимость}

\begin{problem}
    Доказать, что сумма кубов трёх последовательных натуральных чисел 
    всегда делится нацело на $9$.
\end{problem}

\begin{proof}[Решение]~\\
    \begin{quote}
        Обозначим три последовательных натуральных числа через $n - 1$, $n$, 
        $n + 1$ (где $n \ge 2$). Запишем сумму их кубов и раскроем скобки по 
        формулам куба суммы и куба разности:
        \begin{align*}
            (n - 1)^3 &= n^3 - 3n^2 + 3n - 1, \\
            (n + 1)^3 &= n^3 + 3n^2 + 3n + 1.
        \end{align*}
        Складывая все три куба, замечаем, что слагаемые $\mp 3n^2$ и 
        $\mp 1$ взаимно уничтожаются:
        \begin{equation} \label{eq:cubes_sum}
            S(n) = (n - 1)^3 + n^3 + (n + 1)^3 = 3n^3 + 6n \tag{5}
        \end{equation}
        Докажем методом математической индукции, что выражение 
        $S(n) = 3n^3 + 6n$ кратно $9$ при любом натуральном $n$.

        \textbf{1. База индукции.} При $n = 1$ (тройка неотрицательных 
        чисел $0, 1, 2$) имеем $S(1) = 3 \cdot 1^3 + 6 \cdot 1 = 9 \vdots 9$. 
        Для строго натуральных чисел ($n = 2$, тройка $1, 2, 3$) получаем 
        $S(2) = 3 \cdot 2^3 + 6 \cdot 2 = 24 + 12 = 36 \vdots 9$. База верна.

        \textbf{2. Индукционный переход.} Пусть для некоторого $k \ge 1$ 
        число $S(k) = 3k^3 + 6k$ делится на $9$. Рассмотрим выражение 
        $S(k + 1)$ и выделим в нём слагаемое $S(k)$:
        \begin{align*}
            S(k + 1) &= 3(k + 1)^3 + 6(k + 1) = 
            3(k^3 + 3k^2 + 3k + 1) + 6k + 6 \\
            &= (3k^3 + 6k) + 9k^2 + 9k + 9 = S(k) + 9(k^2 + k + 1).
        \end{align*}
        В полученной сумме первое слагаемое $S(k)$ кратно $9$ по 
        индукционному предположению, а второе слагаемое $9(k^2 + k + 1)$ 
        кратно $9$ ввиду наличия сомножителя $9$ при целом $(k^2 + k + 1)$. 
        Следовательно, $S(k + 1) \vdots 9$, что завершает индукцию.
    \end{quote}
\end{proof}

\subsection{Алгоритмические задачи и индукция (Ханойские башни)}

\begin{problem}[\textbf{Задача о Ханойской башне}]
    Даны три стержня. На первом стержне нанизана пирамида из $n$ дисков 
    попарно различного диаметра в порядке убывания диаметра снизу вверх. 
    Требуется перенести всю пирамиду на третий стержень за минимальное 
    число ходов, соблюдая два правила:
    \begin{enumerate}
        \item За один ход разрешается переносить только один верхний диск 
        с одного стержня на другой;
        \item Запрещается класть диск большего диаметра поверх диска 
        меньшего диаметра.
    \end{enumerate}
    Доказать, что минимальное число ходов равно $H_n = 2^n - 1$.
\end{problem}

\begin{proof}[Решение]~\\
    \begin{quote}
        Обозначим через $H_n$ минимальное число ходов, необходимое для 
        переноса пирамиды из $n$ дисков. Докажем формулу $H_n = 2^n - 1$ 
        индукцией по $n$.

        \textbf{1. База индукции ($n = 1$).} Единственный диск переносится 
        с первого стержня на третий за $1$ ход. По формуле $2^1 - 1 = 1$, 
        база выполняется.

        \textbf{2. Индукционный переход ($k \implies k + 1$).} 
        Предположим, что для переноса пирамиды из $k$ дисков минимально 
        требуется $H_k = 2^k - 1$ ходов. Найдём минимальное число ходов 
        для переноса $k + 1$ дисков:
        \begin{enumerate}
            \item Чтобы сдвинуть самый большой ($(k + 1)$-й, нижний) диск 
            с первого стержня на третий, необходимо, чтобы поверх него не 
            было ни одного диска, а третий стержень был полностью свободен. 
            Следовательно, все $k$ меньших дисков должны находиться на 
            втором (вспомогательном) стержне. По индукционному 
            предположению их перенос требует не менее (и в точности при 
            оптимальной игре) $H_k$ ходов.
            \item Перенос самого большого диска с первого стержня на 
            третий занимает ровно $1$ ход.
            \item Теперь требуется перенести пирамиду из $k$ меньших 
            дисков со второго стержня на третий (поверх самого большого 
            диска, который не создаёт ограничений для меньших дисков). 
            По индукционному предположению это требует ещё $H_k$ ходов.
        \end{enumerate}
        Таким образом, выполняется рекуррентное соотношение:
        \begin{equation} \label{eq:hanoi_rec}
            H_{k+1} = H_k + 1 + H_k = 2 H_k + 1 \tag{6}
        \end{equation}
        Подставляя $H_k = 2^k - 1$ в формулу \eqref{eq:hanoi_rec}, получаем:
        \[
            H_{k+1} = 2(2^k - 1) + 1 = 2^{k+1} - 2 + 1 = 2^{k+1} - 1
        \]
        Индукционный переход доказан, значит, $H_n = 2^n - 1$ при всех 
        $n \in \N$.
    \end{quote}
\end{proof}

\subsection{Прямо-обратная индукция (Индукция Коши)}

\begin{definition}
    \textdef{Прямо-обратная индукция} (или \textdef{индукция Коши}) --- 
    модификация метода математической индукции, состоящая из двух этапов:
    \begin{enumerate}
        \item \textdef{Прямой шаг по подпоследовательности:} проверяется 
        база $P(2)$ и доказывается импликация $P(k) \implies P(2k)$, что 
        устанавливает истинность утверждения для всех степеней двойки 
        $n = 2^m$ ($m \in \N$).
        \item \textdef{Обратный шаг (спуск):} доказывается импликация 
        $P(k) \implies P(k - 1)$ для любого $k \ge 2$, позволяющая 
        спуститься от степени двойки $2^m$ к любому предшествующему 
        натуральному числу $n < 2^m$.
    \end{enumerate}
    Совокупность этих двух шагов покрывает всё множество натуральных чисел 
    $\N$.
\end{definition}

\begin{problem}[\textbf{Неравенство Коши о средних}]
    Доказать, что для любых неотрицательных вещественных чисел 
    $x_1, x_2, \dots, x_n \ge 0$ среднее арифметическое не меньше среднего 
    геометрического:
    \begin{equation} \label{eq:am_gm}
        \frac{x_1 + x_2 + \dots + x_n}{n} \ge \sqrt[n]{x_1 x_2 \dots x_n} 
        \tag{7}
    \end{equation}
\end{problem}

\begin{proof}[Решение]~\\
    \begin{quote}
        Применим метод прямо-обратной индукции Коши.

        \textbf{1. Прямой шаг ($k \implies 2k$).} 
        При $n = 2$ неравенство \eqref{eq:am_gm} равносильно 
        неотрицательности полного квадрата:
        \[
            \frac{x_1 + x_2}{2} \ge \sqrt{x_1 x_2} \iff 
            x_1 - 2\sqrt{x_1 x_2} + x_2 \ge 0 \iff 
            \left(\sqrt{x_1} - \sqrt{x_2}\right)^2 \ge 0
        \]
        Пусть неравенство \eqref{eq:am_gm} верно для любых $k$ 
        неотрицательных чисел. Рассмотрим набор из $2k$ чисел 
        $x_1, \dots, x_{2k}$ и разобьём их на две группы по $k$ чисел:
        \[
            A_1 = \frac{x_1 + \dots + x_k}{k}, \qquad
            A_2 = \frac{x_{k+1} + \dots + x_{2k}}{k}.
        \]
        Применяя сначала базу для двух чисел $A_1, A_2$, а затем 
        индукционное предположение для каждой из групп длины $k$, получаем:
        \begin{align*}
            \frac{x_1 + \dots + x_{2k}}{2k} &= \frac{A_1 + A_2}{2} \ge 
            \sqrt{A_1 A_2} \\
            &\ge \sqrt{\sqrt[k]{x_1 \dots x_k} \cdot 
            \sqrt[k]{x_{k+1} \dots x_{2k}}} = 
            \sqrt[2k]{x_1 x_2 \dots x_{2k}}.
        \end{align*}
        Тем самым неравенство \eqref{eq:am_gm} доказано для всех степеней 
        двойки $n = 2, 4, 8, 16, \dots$.

        \textbf{2. Обратный шаг ($k \implies k - 1$).} 
        Предположим, что неравенство \eqref{eq:am_gm} уже доказано для 
        любых $k$ чисел ($k \ge 2$). Возьмём произвольные $k - 1$ 
        неотрицательных чисел $x_1, x_2, \dots, x_{k-1}$ и обозначим их 
        среднее арифметическое через $A$:
        \[
            A = \frac{x_1 + x_2 + \dots + x_{k-1}}{k - 1} \implies 
            x_1 + x_2 + \dots + x_{k-1} = (k - 1)A
        \]
        Дополним этот набор $k$-м числом $x_k = A$ и применим к набору из 
        $k$ чисел $x_1, \dots, x_{k-1}, A$ индукционное предположение:
        \[
            \frac{x_1 + x_2 + \dots + x_{k-1} + A}{k} \ge 
            \sqrt[k]{x_1 x_2 \dots x_{k-1} \cdot A}
        \]
        Преобразуем левую часть неравенства:
        \[
            \frac{(k - 1)A + A}{k} = \frac{k A}{k} = A \implies 
            A \ge \sqrt[k]{x_1 x_2 \dots x_{k-1} \cdot A}
        \]
        Возведём обе части неотрицательного неравенства в степень $k$:
        \[
            A^k \ge x_1 x_2 \dots x_{k-1} \cdot A
        \]
        Если $A = 0$, то $x_1 = \dots = x_{k-1} = 0$, и неравенство 
        очевидно обращается в равенство $0 = 0$. Если же $A > 0$, разделим 
        обе части на $A$ и извлечём корень степени $k - 1$:
        \[
            A^{k-1} \ge x_1 x_2 \dots x_{k-1} \implies 
            A = \frac{x_1 + \dots + x_{k-1}}{k - 1} \ge 
            \sqrt[k-1]{x_1 x_2 \dots x_{k-1}}.
        \]
        Оба шага индукции Коши обоснованы, следовательно, неравенство 
        \eqref{eq:am_gm} верно для всех $n \in \N$.
    \end{quote}
\end{proof}

\subsection{Битовые операции и индукция в информатике}

\begin{definition}
    \textdef{Бинарным (двоичным) представлением} натурального числа 
    $n \in \N$ называется его однозначное разложение по степеням двойки:
    \begin{equation} \label{eq:bin_rep}
        n = \sum_{j=0}^{m} b_j 2^j, \qquad b_j \in \{0, 1\}, \quad b_m = 1, 
        \tag{8}
    \end{equation}
    где индекс старшего значащего разряда равен $m = \lfloor \log_2 n \rfloor$.
\end{definition}

\begin{designation}
    Для двоичного представления натурального числа $n$ введём две 
    числовые характеристики:
    \begin{itemize}
        \item \textdef{$\nu(n)$} --- количество единиц (\textdef{вес 
        Хэмминга}) в двоичной записи числа $n$: $\nu(n) = \sum_{j=0}^m b_j$;
        \item \textdef{$\lambda(n)$} --- номер старшего значащего разряда 
        (количество битов после старшей единицы): 
        $\lambda(n) = \lfloor \log_2 n \rfloor$. Общее количество значащих 
        битов числа $n$ равно $\lambda(n) + 1$.
    \end{itemize}
\end{designation}

\begin{properties}[Рекуррентные свойства функций $\nu(n)$ и $\lambda(n)$] ~
    При приписывании бита справа (умножении на $2$ и возможном прибавлении 
    единицы) выполняются соотношения:
    \begin{align*}
        \nu(2n) &= \nu(n) 
        &&\text{(побитовый сдвиг влево не меняет число единиц)}, \\
        \nu(2n + 1) &= \nu(n) + 1 
        &&\text{(в младший разряд добавляется единица)}, \\
        \lambda(2n) &= \lambda(n) + 1 
        &&\text{(длина записи увеличивается на один разряд)}, \\
        \lambda(2n + 1) &= \lambda(n) + 1 
        &&\text{(длина записи увеличивается на один разряд)}.
    \end{align*}
\end{properties}

\begin{problem}[\textbf{Оптимальное построение числа на калькуляторе}]
    На экране калькулятора изначально написано число $0$. Доступны две 
    операции:
    \begin{enumerate}
        \item Умножить текущее число на $2$ (побитовый сдвиг влево);
        \item Прибавить к текущему числу $1$.
    \end{enumerate}
    Найти минимальное количество операций $\Ops(n)$, необходимое для 
    получения заданного натурального числа $n > 0$.
\end{problem}

\begin{proposition}
    Минимальное количество операций для получения числа $n \in \N$ ($n > 0$) 
    из $0$ выражается формулой:
    \begin{equation} \label{eq:ops_count}
        \Ops(n) = \nu(n) + \lambda(n) \tag{9}
    \end{equation}
\end{proposition}

\begin{proof}[Доказательство]~\\
    \begin{quote}
        Рассмотрим двоичную запись числа $n$ вида \eqref{eq:bin_rep}. 
        Поскольку умножение $0 \cdot 2 = 0$ не меняет начального нуля, 
        первой содержательной операцией обязательно должно быть прибавление 
        единицы ($0 + 1 = 1$), формирующее старший значащий бит $b_m = 1$ 
        (на этот момент затрачена $1$ операция, причём $\nu(1) = 1$, 
        $\lambda(1) = 0$).

        Далее, двигаясь по двоичной записи числа $n$ слева направо от 
        разряда $m - 1$ до разряда $0$:
        \begin{enumerate}
            \item Для перехода к каждому следующему из $\lambda(n)$ 
            разрядов необходимо выполнить операцию умножения на $2$ (сдвиг 
            влево), освобождающую младший бит. Всего потребуется ровно 
            $\lambda(n)$ операций умножения.
            \item Если очередной бит равен $1$, после сдвига выполняется 
            операция $+1$. Поскольку старшая единица уже была получена на 
            первом шаге, потребуется ещё ровно $\nu(n) - 1$ операций 
            прибавления единицы.
        \end{enumerate}
        Суммируя число операций, получаем:
        \[
            \Ops(n) = 1 + \lambda(n) + \bigl(\nu(n) - 1\bigr) = 
            \nu(n) + \lambda(n)
        \]
        Строгое доказательство оптимальности проводится сильной индукцией 
        по $n$ с использованием рекуррентных свойств функций $\nu$ и 
        $\lambda$: переход к чётному числу $2k$ выгоднее делать удвоением 
        из $k$, а к нечётному $2k + 1$ --- единственным возможным шагом 
        $+1$ из чётного числа $2k$.
    \end{quote}
\end{proof}

\begin{problem}[\textbf{Сумма всех циклических сдвигов числа}]
    Рассмотрим $n$-битовое число $x = \sum_{j=0}^{n-1} b_j 2^j$, где 
    $b_j \in \{0, 1\}$. Обозначим через $(x \lll i)$ результат 
    циклического побитового сдвига числа $x$ влево на $i$ разрядов. 
    Доказать, что сумма всех $n$ циклических сдвигов числа $x$ равна:
    \begin{equation} \label{eq:cyc_shift_sum}
        \sum_{i=0}^{n-1} (x \lll i) = (2^n - 1) \cdot \nu(x) \tag{10}
    \end{equation}
\end{problem}

\begin{proof}[Решение]~\\
    \begin{quote}
        При циклическом сдвиге $n$-битового слова влево на $i$ позиций бит 
        $b_j$, стоявший в $j$-м разряде, перемещается в разряд с номером 
        $(j + i) \bmod n$. Следовательно, числовое значение сдвинутого 
        слова записывается в виде:
        \[
            (x \lll i) = \sum_{j=0}^{n-1} b_j 2^{(j+i) \bmod n}
        \]
        Просуммируем это равенство по всем величинам сдвига 
        $i = 0, 1, \dots, n - 1$ и поменяем порядок конечного суммирования:
        \[
            S = \sum_{i=0}^{n-1} (x \lll i) = 
            \sum_{i=0}^{n-1} \sum_{j=0}^{n-1} b_j 2^{(j+i) \bmod n} = 
            \sum_{j=0}^{n-1} b_j \left( \sum_{i=0}^{n-1} 2^{(j+i) \bmod n} \right)
        \]
        Проанализируем внутреннюю сумму по $i$ при фиксированном индексе 
        $j$. Когда переменная $i$ пробегает полную систему вычетов 
        $0, 1, \dots, n - 1$, величина $(j + i) \bmod n$ также пробегает 
        все значения от $0$ до $n - 1$ ровно по одному разу. По формуле 
        суммы геометрической прогрессии находим:
        \[
            \sum_{i=0}^{n-1} 2^{(j+i) \bmod n} = 
            \sum_{k=0}^{n-1} 2^k = 2^0 + 2^1 + \dots + 2^{n-1} = 2^n - 1
        \]
        Заметим, что значение внутренней суммы не зависит от $j$. Вынося 
        множитель $(2^n - 1)$ за знак внешней суммы и учитывая, что 
        $\sum_{j=0}^{n-1} b_j = \nu(x)$, окончательно получаем:
        \[
            S = (2^n - 1) \sum_{j=0}^{n-1} b_j = (2^n - 1) \cdot \nu(x)
        \]
        Тождество \eqref{eq:cyc_shift_sum} полностью доказано.
    \end{quote}
\end{proof}